% Regression: p3_corner — imported current-behavior fixture from the handoff corpus.
% Formula: patch P = 3x4 window; erased set E = {(1,1),(1,2)} u {(3,3),(3,4)}
% p3_corner — arXiv:1804.04964 (normal PEPS FT), Corollary 8 sketch, p. 24.
% Corner-erasure sequence: compare the first and last sliding windows on a
% 3x4 open patch; adding two-two tensors in the upper-left and lower-right
% corners and inverting the resulting regions ERASES those corners, and on
% the erased patch the two window insertions are equal.
%
% Formulae:
%   patch P = 3x4 window; erased set E = {(1,1),(1,2)} u {(3,3),(3,4)}
%   (the "two-two tensors" the paper removes by inversion);
%   on P \ E:   O_{W1}|psi>  =  O_{W4}|psi>,
%   W1 = (1-2, 3-4) (window at the north-east), W4 = (2-3, 1-2) (window at
%   the south-west), the distinguished edge e = <(2,2),(2,3)> fixed.
% The top row shows the erasure step itself as a tenkzcd arrow.
%
% Divergence noted: the paper's white-out erasure cuts bonds midway,
% leaving dangling half-bond stubs on the surviving neighbours (the cut
% virtual indices opened by the erasure); tenkz's \tnsite[removed] deletes
% the site WITH its bonds and draws no cut stubs toward the erased cells.
\documentclass[varwidth=26cm, border=6pt]{standalone}
\usepackage{amsmath, amssymb}
\usepackage{tenkz}
\begin{document}
% ---- the erasure step: full patch -> corner-erased patch ----
\begin{tenkzcd}[column sep=huge]
  \begin{tenkzlattice}[rows=3, cols=4, boundary legs, compact]
    \tnregion[slot=selected, outline]{(1-2,3-4)}
    \tnedge{(2,2)-(2,3)}
  \end{tenkzlattice}
  \arrow[r, "\text{erase corners}"]
  &
  \begin{tenkzlattice}[rows=3, cols=4, boundary legs, compact]
    \tnsite[removed]{(1,1)} \tnsite[removed]{(1,2)}
    \tnsite[removed]{(3,3)} \tnsite[removed]{(3,4)}
    \tnregion[slot=selected, outline]{(1-2,3-4)}
    \tnedge{(2,2)-(2,3)}
  \end{tenkzlattice}
\end{tenkzcd}

\bigskip
% ---- the resulting equality on the erased patch ----
$
\begin{tenkzlattice}[rows=3, cols=4, boundary legs, compact]
  \tnsite[removed]{(1,1)} \tnsite[removed]{(1,2)}
  \tnsite[removed]{(3,3)} \tnsite[removed]{(3,4)}
  \tnregion[slot=selected, outline]{(1-2,3-4)}
  \tnedge{(2,2)-(2,3)}
\end{tenkzlattice}
\;=\;
\begin{tenkzlattice}[rows=3, cols=4, boundary legs, compact]
  \tnsite[removed]{(1,1)} \tnsite[removed]{(1,2)}
  \tnsite[removed]{(3,3)} \tnsite[removed]{(3,4)}
  \tnregion[slot=selected, outline]{(2-3,1-2)}
  \tnedge{(2,2)-(2,3)}
\end{tenkzlattice}
$
\end{document}
